\documentclass[a4paper,10pt]{article}
\usepackage[utf8]{inputenc}
\usepackage[T1]{fontenc}
\usepackage[spanish]{babel}
\usepackage{geometry}
\geometry{top=0.5in, bottom=0.4in, left=0.55in, right=0.55in}
\usepackage{enumitem}
\setlist[itemize]{noitemsep, topsep=2pt, leftmargin=0.25in, itemsep=1pt, label=\textbullet}
\usepackage[hidelinks]{hyperref}
\usepackage{titlesec}
\usepackage{helvet}
\renewcommand{\familydefault}{\sfdefault}

% --- Harvard-style section: centered, bold, with rule below ---
\titleformat{\section}{\normalsize\bfseries\centering}{}{0em}{\MakeUppercase}[\vspace{2pt}\titlerule]
\titlespacing*{\section}{0pt}{8pt}{5pt}

\setlength{\parskip}{0pt}
\setlength{\parsep}{0pt}
\setlength{\parindent}{0pt}

% Spacing helpers
\newcommand{\entryspace}{\vspace{5pt}}
\newcommand{\subentryspace}{\vspace{2pt}}

\begin{document}

% ===================== HEADER =====================
\begin{center}
    {\Large\bfseries Jonathan Sánchez Pesantes} \\[4pt]
    {\normalsize +569 6758 5503 \enspace\textbullet\enspace \href{mailto:jonathan.sanchez.ep@gmail.com}{jonathan.sanchez.ep@gmail.com} \enspace\textbullet\enspace \href{https://www.linkedin.com/in/jonasanchez}{linkedin.com/in/jonasanchez}} \enspace\textbullet\enspace \href{https://github.com/jsanchez-ds}{github.com/jsanchez-ds}
\end{center}

% ===================== SUMMARY =====================
\section*{Resumen Profesional}
Ingeniero Civil Industrial (U. de Chile, distinción máxima) con interés en ciencia de datos y modelamiento predictivo. Experiencia en automatización de procesos con Python, control de gestión y desarrollo de modelos de forecasting con machine learning. Manejo de Python, R y SQL. Busco roles de Ciencia de Datos donde pueda seguir desarrollando habilidades en ML aplicado a problemas de negocio.

% ===================== EXPERIENCE =====================
\section*{Experiencia Profesional}

\textbf{ClaroVTR} \hfill Santiago, Chile \\
\textbf{Ingeniero en Eficiencias} -- Subgerencia de Gestión Operacional \hfill Marzo 2025 -- Presente
\begin{itemize}
    \item Diseñé modelo de predicción de consumo de datos (XGBoost + Prophet) para un servicio con $\sim$500 líneas, incluyendo feature engineering y validación temporal.
    \item Automaticé procesos operativos con Python: reportes, extracción de datos desde API REST Starlink (OAuth2), conciliación con SAP y SQL.
    \item Gestioné iniciativas de eficiencia operacional enfocadas en reducción de costos y optimización tarifaria.
    \item Llevé procesos de control de gestión y validación de facturación a nivel corporativo.
    \item Construí dashboards en Power BI para seguimiento de KPIs operacionales.
\end{itemize}

\entryspace

\textbf{Keirón} \hfill Santiago, Chile \\
\textbf{Analyst -- Customer Success} \hfill 2024
\begin{itemize}
    \item Diseñé e implementé el \textit{Customer Health Index} (CHI), modelo de scoring multivariable para predicción de riesgo de fuga en segmentos Enterprise.
    \item Analicé patrones de uso para detección temprana de clientes en riesgo, integrando múltiples señales de adopción y actividad.
    \item Desarrollé dashboards en Redash para métricas clave de adopción y uso de plataforma.
\end{itemize}

\entryspace

\textbf{Agrosuper} \hfill Santiago, Chile \\
\textbf{Proyecto de Título -- Ingeniería en Planificación} \hfill 2024
\begin{itemize}
    \item Implementé scripts en Python para limpieza y validación masiva de datos de importaciones, detectando inconsistencias históricas en registros de comercio exterior.
    \item Reemplacé cruces manuales en Excel por procesos automatizados de conciliación entre SAP y bases SQL, reduciendo significativamente tiempos de procesamiento y mejorando trazabilidad de inventarios.
\end{itemize}

\entryspace

\textbf{Cencosud Scotiabank} \hfill Santiago, Chile \\
\textbf{Práctica Profesional -- Analista MIS} \hfill 2023
\begin{itemize}
    \item Optimicé consultas SQL (Query Tuning) sobre grandes volúmenes de datos financieros e integré múltiples fuentes mediante procesos ETL para el Panel de Información Institucional.
    \item Generé reportes ejecutivos de KPIs para la alta gerencia en Chile y Canadá.
\end{itemize}

% ===================== EDUCATION =====================
\section*{Formación Académica}

\textbf{Pontificia Universidad Católica de Chile} \hfill Santiago, Chile \\
Diplomado en Big Data y Ciencia de Datos (296 horas) \hfill 2026 \\
Infraestructura de Datos, Minería de Datos, Visualización, Programación en R y Aprendizaje de Máquina.

\entryspace

\textbf{Universidad Adolfo Ibáñez} \hfill Santiago, Chile \\
Gestión y Mejora de Procesos -- BPM (96 horas) \hfill 2025 \\
Mejora continua, diseño de KPIs y evaluación de calidad de datos.

\entryspace

\textbf{Universidad de Chile} \hfill Santiago, Chile \\
Ingeniero Civil Industrial, Licenciado en Ciencias de la Ingeniería mención Industrial \hfill 2019 -- 2024 \\
Titulado con distinción máxima. Estudiante Destacado 2023 -- Top $\sim$6\% de la generación.

% ===================== PROJECTS =====================
\section*{Proyectos Analíticos}

\textbf{\href{https://bank-marketing-analysis-jsanchez.streamlit.app/}{Análisis de Campañas Bancarias}} -- \textit{Demo interactiva en Streamlit} \hfill 2024 \\
Predicción de suscripción a depósito a plazo (dataset UCI, 45K registros). Random Forest con ROC-AUC 0.80 sobre hold-out (sin \texttt{duration} para evitar data leakage). Diagnóstico y corrección de fuga SMOTE-en-CV.

\entryspace

\textbf{\href{https://jsanchez-ds.github.io/medical-diagnosis-classification/}{Clasificación de Diagnóstico Médico con ML}} \hfill 2024 \\
Clasificación benigno/maligno sobre dataset Wisconsin (569 registros, 30 variables). SVM con GridSearchCV alcanza 97.6\% de accuracy y AUC 0.99. Manejo de desbalance con Under/Over-Sampling.

\entryspace

\textbf{\href{https://jsanchez-ds.github.io/gender-income-gap/}{Brecha de Género en Ingresos de Pequeños Comercios}} \hfill 2023 \\
Regresión con efectos fijos (\texttt{fixest}) sobre $\sim$5.000 comercios. Selección de variables con LASSO/Ridge. MARS seleccionado por menor error out-of-sample.

\entryspace

\textbf{\href{https://jsanchez-ds.github.io/credit-choice-experiment/}{Experimento de Elección de Crédito Financiero}} \hfill 2023 \\
Discrete choice modeling con 4 condiciones experimentales. Logit condicional, Mixed Logit y comparación ML (RF, SVM, CART, KNN).

% ===================== SKILLS =====================
\section*{Habilidades Técnicas e Idiomas}

\textbf{Lenguajes:} Python (Pandas, NumPy, scikit-learn), R (tidyverse, ggplot2), SQL \\
\textbf{ML \& Stats:} XGBoost, Random Forest, SVM, LASSO/Ridge, Logit, Prophet, regresión con efectos fijos \\
\textbf{Técnicas:} Feature engineering, validación cruzada, manejo de desbalance, selección de variables \\
\textbf{Herramientas:} Power BI, PySpark, APIs REST, Git, SAP \\
\textbf{Bases de Datos:} PostgreSQL, SQL Server \\
\textbf{Idiomas:} Español (Nativo), Inglés (Intermedio Profesional)

\end{document}
